\begin{ledgroupsized}[r]{120mm}
\footnotesize
\pstart
\noindent\textbf{Überlieferung:}
\pend
\end{ledgroupsized}
\begin{ledgroupsized}[r]{114mm}
\footnotesize
\pstart
\parindent -6mm
\makebox[6mm][l]{\textit{L}}%
Konzept:
LH~XXXVII~3 Bl.~82\textendash83. 
Ein Bogen~2\textsuperscript{o};
\href{\Scan{LH_37_03_083wz}}%
{Wasserzeichen} in Bl.~83.
Drei Seiten auf Bl.~82~r\textsuperscript{o} bis Bl.~83~v\textsuperscript{o};
Textfolge: Bl.~82~r\textsuperscript{o}, obere Hälfte von Bl.~82~v\textsuperscript{o}, Bl.~83~r\textsuperscript{o}, untere Hälfte von Bl.~82~v\textsuperscript{o}; 
Bl.~83~v\textsuperscript{o} leer. 
Bei der Anfertigung der Zeichnung \bracks{\textit{Fig.~2}} auf Bl.~\href{\Scan{LH_37_03_082v}}{82~v\textsuperscript{o}} (hier S.~\pageref{Fig:37_03_082-083:Fig_1}) hat Leibniz sich an den durch das Blatt durchscheinenden Konturen der \bracks{\textit{Fig.~1}} auf der Rückseite (Bl.~\href{\Scan{LH_37_03_082r}}{82~r\textsuperscript{o}}, S.~\pageref{Fig:37_03_082-083:Fig_2}) orientiert, die damit als exakte Vorlage für \bracks{\textit{Fig.~2}} diente.
\pend
\end{ledgroupsized}
\selectlanguage{latin}
\vspace{8mm}
\pstart
\noindent
\href{\Scan{LH_37_03_082r}}%
{\lbrack82~r\textsuperscript{o}\rbrack}\ 
\pend
\nobreak	%
\begin{edfigure}[.5em][1em]
\includegraphics[scale=1]{%
gesamttex/edit_VIII-4/images/LH_37_03_082-083_d1_082r.pdf}
\edcaption[0.2em]{\lbrack\textit{Fig.~1, tlw. Blindzeichnung}\rbrack%
\pagelabel{Fig:37_03_082-083:Fig_1}}
\end{edfigure}
\pstart
\noindent
\textls{Centrum gravitatis}\indexS{centrum gravitatis} est punctum ex quo suspensum\indexS{corpus ex centro gravitatis suspensum}
\edtext{corpus, non movetur vi gravitatis.\indexS{gravitas}}{\lemma{corpus,}\Bfootnote{%
\textit{(1)}~vi gravitatis %
\textit{(2)}~non movetur vi gravitatis. %
\textit{(a)}~\bars{\bars{Itaque si \textit{streicht Hrsg.}} figurae situs \textit{erg.}} %
\textit{(b)}~\textls{Centrum}~\textit{L}}}
\pend
\pstart
\textls{Centrum figurae}\indexS{centrum figurae} est punctum per quod transiens recta vel planum quodlibet figuram secat in duas partes similes\bracks{,} aequales, et ex adverso respondentes 
\edtext{seu similiter positas.}{\lemma{}\Bfootnote{seu similiter positas \textit{erg.~L}}} 
\textls{Figuram}\indexS{figura Amphidextra} autem centrum habentem\indexS{figura centrum habentem} possis vocare \textls{Amphidextram} 
\edtext{ut Linea Recta,}{\lemma{ut}\Bfootnote{%
\textit{(1)}~Circulus, %
\textit{(2)}~Linea Recta,~\textit{L}}} 
Peripheria circuli, Circulus, Ellipsis. Sphaera\bracks{,} Sphaeroides, Cylinder.
\pend
\pstart
\textls{Omnis figura amphidextra\indexS{figura amphidextra} habet centrum gravitatis\indexS{centrum gravitatis} idem cum centro figurae.\indexS{centrum figurae}}
Hoc manifestum est, quia a qualibet recta vel 
\edtext{plano per}{\lemma{plano}\Bfootnote{%
\textit{(1)}~ita secatur, %
\textit{(2)}~per~\textit{L}}}
centrum gravitatis\indexS{centrum gravitatis} ita secatur, ut omnium par sit ratio ab utraque 
\edlabel{37_03_082-083_1a}%
\edtext{}{%
\xxref{37_03_082-083_1a}{37_03_082-083_1b}%
\lemma{parte.}%
\Bfootnote{%
\lbrack/\rbrack\
\textit{(1)}~Figura %
\textit{(2)}~\textls{Duarum}~\textit{L}}}%
parte. 
\pend
\pstart
\textls{Duarum\edlabel{37_03_082-083_1b} figurarum liberarum\indexS{figura libera} vel ex centro gravitatis\indexS{centrum gravitatis} suspensarum\indexS{figura ex centro gravitatis suspensa} pondera\indexS{pondus} sunt inter se, ut earum areae}. 
\pend
\pstart
Si corporis situs utcunque mutetur, manente situ\indexS{situs centri gravitatis} centri gravitatis\indexS{centrum gravitatis} nihil variatur quoad gravitatem.
\pend
\vspace{.5\baselineskip}
\pstart
\centering
\edtext{\textls{Theorema aequiponderantium}.\indexS{Theorema aequiponderantium}}{\lemma{\textls{Theorema aequiponderantium}}\Cfootnote{Doppelt unterstrichen.}}
\pend
\nobreak
\pstart
Hoc ita investigavi, methodo analytica.\indexS{methodus analytica}
\pend
\pstart
Duo 
\edtext{corpora\indexS{corpora ex extremis rectae suspensa} \textit{ACB} et \textit{DEF}, suspensa}{\lemma{}\Bfootnote{%
corpora
\bars{%
\textit{(1)}~\textit{AB} et \textit{DF} %
\textit{(2)}~\textit{ACB} et \textit{DEF}, \textit{erg.}} %
\textit{(1)}~ex eadem libra debent esse suspensa per centrum gravitatis %
\textit{(2)}~suspensa~\textit{L}}}
sunt centris gravitatis\indexS{centrum gravitatis} \textit{C} et \textit{E} ex rectae \textit{CE} 
\edtext{extremis, et sit}{\lemma{extremis,}\Bfootnote{%
\textit{(1)}~quaeritur %
\textit{(2)}~et sit~\textit{L}}}
punctum aequilibrii\indexS{punctum aequilibrii}\indexS{aequilibrium} \textit{G}. Quaeritur ratio \textit{CG} ad \textit{GE}. 
\pend
\pstart
Haec 
\edtext{ut investigetur,}{\lemma{ut}\Bfootnote{%
\textit{(1)}~facilius investigetur, %
\textit{(2)}~investigetur,~\textit{L}}}
considerandum est, corpora \textit{AB}, \textit{DF} varias habere posse 
\edtext{figuras, manente}{\lemma{figuras,}\Bfootnote{%
\textit{(1)}~easque nihil %
\textit{(2)}~manente~\textit{L}}}
eodem pondere,\indexS{pondus corporis}
et assumta eorum figura\indexS{figura corporis} 
quacunque varios etiam situs eorum intelligi posse, 
manente semper eodem centri gravitatis situ,\indexS{situs centri gravitatis}\indexS{centrum gravitatis} 
adeoque et manente aequilibrio.\indexS{aequilibrium} Et cum fieri possit ut una figura alia aptior sit ad hanc 
\edtext{indagationem; et variarum}{\lemma{indagationem;}\Bfootnote{%
\textit{(1)}~hinc %
\textit{(2)}~et %
\bars{fortasse \textit{gestr.}} variarum~\textit{L}}}
figurarum, variorumque situum inter se collatio institui possit, hinc ordine varios status examinabimus. 
\pend
\pstart
Simplicissimus omnium est, ut totum corpus \textit{AB}, vel \textit{DF}, intelligatur collectum;\indexS{pondus corporis in punctum collectum} illud in punctum \textit{C}, hoc in punctum \textit{E}. 
\pend
\pstart
Sed 
\edtext{cum hoc modo}{\lemma{cum}\Bfootnote{%
\textit{(1)}~huic %
\textit{(2)}~hoc modo~\textit{L}}}
in partes dividi non possint 
\edtext{corpora, hinc utilis non}{\lemma{corpora}\Bfootnote{%
\textit{(1)}~nec %
\textit{(2)}~hinc utilis non~\textit{L}}}
est haec consideratio. Proximum ex 
\edtext{simplicissimis hoc}{\lemma{simplicissimis}\Bfootnote{%
\textit{(1)}~hoc %
\textit{(2)}~hoc~\textit{L}}}
est, ut ponamus utrumque corpus extensum\indexS{corpus extensum} in lineam rectam \textit{ACB} vel \textit{DEF}. Quarum situs maxime simplex\indexS{situs corporis maxime simplex} et ille est,
\edtext{ut ad rectam}{\lemma{ut}\Bfootnote{%
\textit{(1)}~sint para %
\textit{(2)}~ad %
\textit{(a)}~centrum aequi %
\textit{(b)}~rectam aequilib %
\textit{(c)}~rectam~\textit{L}}}
\textit{CE} angulum faciant rectum, adeoque sint parallelae inter 
\edlabel{37_03_082-083_2a}%
\edtext{}{%
\xxref{37_03_082-083_2a}{37_03_082-083_2b}%
\lemma{se.}%
\Bfootnote{%
\textit{(1)}~Per punctum \textit{G} quaesitum ducatur recta \textit{PGQ} %
\textit{(a)}~illis parallela %
\textit{(b)}~quae erit alius aequilibrii axis %
\textit{(2)}~Porro~\textit{L}}}%
se. 
\pend
\pstart
Porro\edlabel{37_03_082-083_2b}
puncti quaesiti \textit{G} ea est natura\bracks{,} 
\edtext{seu punctum tale a nobis quaeritur,}{\lemma{}\Bfootnote{seu punctum tale a nobis quaeritur \textit{erg.~L}}}
ut quaecunque recta per ipsum transiens sit axis aequilibrii\indexS{axis aequilibrii}\indexS{aequilibrium}
duorum 
\edtext{corporum \textit{AB}, \textit{DF} ita}{\lemma{corporum \textit{AB}, \textit{DF}}\Bfootnote{%
\textit{(1)}~ita ut corpora %
\textit{(2)}~. Et cum variae %
\textit{(3)}~ita~\textit{L}}}
ut si ipsa sustineat totam libram \textit{CE} cum corporibus annexis, mansurum sit aequilibrium.\indexS{aequilibrium} %
\pend
\pstart
\edlabel{37_03_082-083_6a}%
Ex pluribus autem qui duci 
\edtext{\bracks{possunt}}{\lemma{}\Bfootnote{possunt \textit{erg.~Hrsg.}}}
\href{\Scan{LH_37_03_082v}}%
{\lbrack82~v\textsuperscript{o}\rbrack}
aequilibrii\indexS{aequilibrium}
\edtext{axibus\indexS{axis aequilibrii} ducamus}{%
\lemma{axibus}%
\Bfootnote{%
\textit{(1)}~ducatur
\textit{(2)}~ducamus~\textit{L}}}
\textit{PQ} parallelum corporibus, seu ipsi librae \textit{CE} perpendicularem 
\edtext{in \textit{G}. Jam quia rursus ponimus eam}{%
\lemma{in \textit{G}.}%
\Bfootnote{%
\textit{(1)}~Ubi rursus manifestum est %
\textit{(2)}~Jam quia rursus %
\textit{(a)}~manifestum est %
\textit{(aa)}~manente eadem distantia ab axe aequilibrii %
\textit{(b)}~ponimus eam~\textit{L}}}
esse naturam puncti \textit{G} quaesiti seu rectae \textit{PGQ} per ipsum transeuntis, ut manente centri 
\edtext{gravitatis\indexS{centrum gravitatis} cujusque}{%
\lemma{gravitatis cujusque}%
\Bfootnote{%
\textit{(1)}~eadem distantia
\textit{(2)}~cujusque~\textit{L}}}
corporis eadem distantia\indexS{distantia centri gravitatis ab axe aequilibrii} ab axe aequilibrii\indexS{axis aequilibrii} idem maneat aequilibrium,\indexS{aequilibrium} 
\edtext{hinc patet}{\lemma{hinc}\Bfootnote{%
\textit{(1)}~servato %
\textit{(2)}~patet~\textit{L}}}
\textit{A\,C\,B} transferri posse in (\textit{A})\,(\textit{C})\,(\textit{B}) in eadem recta \textit{A}(\textit{A})\textit{C}(\textit{C})\textit{B}(\textit{B}).%
\edlabel{37_03_082-083_6b}%
\pend
\begin{edfigure}[1.2em][1em]
\includegraphics[scale=1]{%
gesamttex/edit_VIII-4/images/LH_37_03_082-083_d2_082v.pdf}
\edcaption[.25em]{\edtext{\lbrack\textit{Fig.~2, tlw. Blindzeichnung}\rbrack\pagelabel{Fig:37_03_082-083:Fig_2}}{%
\lemma{\lbrack\textit{Fig.~2}\rbrack}%
\Cfootnote[nonum]{Bei der Anfertigung der Zeichnung hat Leibniz sich an den durch das Blatt durchscheinenden Konturen der \lbrack\textit{Fig.~1}\rbrack\ auf der Rückseite (Bl.~\href{\Scan{LH_37_03_082r}}{82~r\textsuperscript{o}}) orientiert, die damit als exakte Vorlage für \lbrack\textit{Fig.~2}\rbrack\ diente.}}}%
\end{edfigure}
\pstart
Cumque varie assumi possit distantia \textit{A}(\textit{A}) vel \textit{B}(\textit{B}), eligamus eam quae 
\edtext{nobis ad}{\lemma{}\Bfootnote{nobis \bars{aliquid \textit{gestr.}} ad~\textit{L}}}
contemplationem utilis esse possit, itaque eousque moveamus \textit{AB} in (\textit{A})(\textit{B})  donec juncta \textit{D}(\textit{B}) secet ambas rectas \textit{PQ}, et 
\edtext{(\textit{C})\lbrack\textit{E}\rbrack}{%
\lemma{}%
\Bfootnote{(\textit{C})\textit{R} \textit{L ändert Hrsg.}}}
(ipsi \textit{CE} aequalem et parallelam) in eodem puncto (\textit{G}). Hic autem fieri posse patet, nam 
\edtext{si angulus}{\lemma{}\Bfootnote{si \bars{totus \textit{gestr.}} angulus~\textit{L}}}
\textit{ABGS} rectae \textit{CE} affixus, et cum ea descendens intelligatur, perveniet \textit{AB} in (\textit{A})(\textit{B}) et \textit{CE} in (\textit{C})(\textit{E}) et \textit{G} in (\textit{G}) et \textit{S} in (\textit{S}) adeoque tamdiu descendet \textit{S} per rectam \textit{SD}, donec incidat (\textit{S}) in 
\edlabel{37_03_082-083_3a}%
\edtext{}{\xxref{37_03_082-083_3a}{37_03_082-083_3b}%
\lemma{ipsam \textit{D}.}%
\Bfootnote{%
\textit{(1)}~Patet \textit{B} \textit{G} \textit{D} fore etiam axem aequilibri. Adeoque 
\textit{(a)}~distantia \textit{E} centri gravitatis ipsius \textit{DF}, \bars{nempe \textit{EM} \textit{erg.}} 
\textit{(b)}~\textit{EM} distantia \textit{E} centri gravitatis ipsius \textit{DF}, ab axe 
aequilibrii \textit{DB} erit ad (\textit{C})\textit{N} distantiam (\textit{C}) centri gravitatis ipsius (\textit{A})(\textit{B}) ut %
\textit{(aa)}~\textit{FG} ad %
\textit{(bb)}~(\textit{E})(\textit{G}) ad (\textit{G})(\textit{C}) seu ut \textit{EG} ad \textit{GC}. Atque ita ad aliquam aequationem sive cognitionem pervenimus quae ad rem definiendam sufficit. 
\lbrack/\rbrack\ Nimirum ob triangula similia \textit{EMD} et (\textit{C})\textit{N}(\textit{B}) erit (\textit{C})(\textit{B}) ad (\textit{C})\textit{N} ut
\href{\Scan{LH_37_03_083r}}%
{\lbrack83~r\textsuperscript{o}\rbrack}  
\textit{(2)}~\bars{Sit \textit{streicht Hrsg.}} pa
\textit{(3)}~Sit~\textit{L}}}%
ipsam \textit{D}.
\pend 
\clearpage
\pstart
\noindent
\href{\Scan{LH_37_03_083r}}%
{\lbrack83~r\textsuperscript{o}\rbrack}%
\pend
\nobreak
\begin{edfigure}[1em][1em]
\hfill
\includegraphics[scale=1]{%
gesamttex/edit_VIII-4/images/LH_37_03_082-083_d3_083r.pdf}
\hfill%
\includegraphics[scale=1]{%
gesamttex/edit_VIII-4/images/LH_37_03_082-083_d4_083r.pdf}
\hfill
\edcaption{%
\hspace*{5mm}%
\lbrack\textit{Fig.~3}\rbrack%
\hspace*{30mm}%
\lbrack\textit{Fig.~4, tlw. Blindzeichnung}\rbrack}
\end{edfigure}
\pstart 
Sit\edlabel{37_03_082-083_3b}
parallelogrammum\indexS{parallelogrammum} \textit{DVBF}, in cujus duobus lateribus \textit{DF}, \textit{VB} sumantur duo quaecunque puncta \textit{E}, \textit{C} 
velut poli circa quos parallelogrammum\indexS{parallelogrammum circa polos duos mobile} 
moveatur. Jungatur \textit{DB} secans ipsam \textit{EC} in \textit{T}. Ajo ipsam (\textit{B})(\textit{D}) eandem rectam immobilem \textit{EC} secare videri in eodem puncto \textit{T}, seu punctum quoque \textit{T} esse immobile. Hoc an verum sit, quaerenda demonstratio 
\edlabel{37_03_082-083_5a}%
\edtext{}{%
\xxref{37_03_082-083_5a}{37_03_082-083_5b}%
\lemma{est.}%
\Bfootnote{%
\lbrack/\rbrack\
\textit{(1)}~Nimirum hoc esse deberet \bars{ut \textit{streicht Hrsg.}} duorum circulorum %
\textit{(a)}~centra \textit{E} \textit{F} %
\textit{(b)}~diametri sint %
\textit{(c)}~radii sint paralleli \textit{DEF}, \textit{ACB}. %
\textit{(2)}~Sed ea~\textit{L}}}%
est.
\pend
\pstart
Sed ea\edlabel{37_03_082-083_5b}
ut video est facilis. Nam semper \textit{DE} ad \textit{ET}, ut \textit{BV} ad \textit{CT}, quae ratio cum semper sit eadem, etiam \textit{ET} ad \textit{TC} semper erit eadem, et 
\edtext{cum \textit{EC}}{\lemma{cum}\Bfootnote{%
\textit{(1)}~\textit{ETC} semper sit %
\textit{(2)}~\textit{EC}~\textit{L}}}
sit semper idem, etiam punctum \textit{T} erit semper idem. Non ergo effici potest, ut recta \textit{BD} incidat in \textit{G}, sed incidet in \textit{T} semper. Quia 
\edtext{autem patet}{\lemma{}\Bfootnote{%
autem \bars{hinc \textit{gestr.}} patet~\textit{L}}}
\textit{ET} esse aequ. \textit{CG}, tunc ut in \textit{G} incidat linea jungens duo puncta \textit{D}, \textit{B} efficiemus fortasse sumendo \textit{DF} et \textit{AB} non parallelas sed alteram alteri 
\edtext{si opus productae}{\lemma{}\Bfootnote{si opus productae \textit{erg.~L}}} 
perpendicularem. Compleatur parallelogrammum \textit{ALMB} ita ut sit \textit{AL} vel \textit{MB} parall. et aequ. \textit{EC} et \textit{LM} parall. et aequ. \textit{AB}. 
\pend 
\begin{edfigure}[1em][1em]
\hspace*{-5em}%
\includegraphics[scale=1]{%
gesamttex/edit_VIII-4/images/LH_37_03_082-083_d5_083r.pdf}
\edcaption{\lbrack\textit{Fig.~5, tlw. Blindzeichnung}\rbrack}
\end{edfigure}
\begin{edfigure}[1em][1em]
\includegraphics[scale=1]{%
gesamttex/edit_VIII-4/images/LH_37_03_082-083_d6_083r.pdf}
\edcaption{%
\edtext{\lbrack\textit{Fig.~6}\rbrack}{%
\lemma{\lbrack\textit{Fig.~6}\rbrack}%
\Cfootnote[nonum]{Die von Leibniz gestrichene Strecke (\textit{A})\textit{H} ergänzt Hrsg.}}
}
\end{edfigure}
\pstart
\edlabel{37_03_082-083_inventio-a}%
Sed nunc tandem veram et pulcherrimam et simul
\edtext{analyticam theorematis\indexS{inventio analytica theorematis}}{\lemma{analyticam}\Bfootnote{%
\textit{(1)}~problematis com %
\textit{(2)}~theorematis~\textit{L}}}
inventionem reperi, quae ita habet.%
\edlabel{37_03_082-083_inventio-b}
\pend
\pstart
Sint duo 
\edtext{pondera \textit{C}, \textit{E}}{%
\lemma{pondera}%
\Bfootnote{%
\textit{(1)}~,~\textit{ACB}, \textit{DEF} ex n
\textit{(2)}~\textit{E}, \textit{C}
\textit{(3)}~\textit{C}, \textit{E}~\textit{L}}}
repraesentata per rectas (vel cylindros)\indexS{cylinder centrum gravitatis habens} \textit{ACB} et
\edtext{\textit{DEF} centra gravitatis\indexS{centrum gravitatis} \textit{C}, \textit{E} habentes. Nihil}{%
\lemma{}%
\Bfootnote{%
\textit{DEF} %
\bars{centra gravitatis \textit{C}, \textit{E} habentes~\textit{erg.}}
\bars{centra gravitatis \textit{C}. \textit{E} habentes \textit{erg., streicht Hrsg.}}.~Nihil~\textit{L}}}
enim refert quomodo figura ponderis\indexS{figura ponderis} mutetur modo idem maneat situs centri\indexS{situs centri gravitatis} 
\edtext{gravitatis\indexS{centrum gravitatis}}{\lemma{}\Bfootnote{gravitatis \textit{erg.~L}}} 
id est puncti per quod transiens quodlibet planum figuram secat
in partes aequiponderantes,\indexS{partes aequiponderantes figurae} quale punctum in figuris amphidextris\indexS{figura amphidextra} 
dari manifestum est, in quibus coincidit cum centro figurae.\indexS{centrum figurae} 
Jungantur eorum centra \textit{E}, \textit{C} recta \textit{EC} velut libra cujus fulcrum\indexS{fulcrum librae} sit alicubi in \textit{G} puncto aequilibrii,\indexS{punctum aequilibrii} ita ut ex 
\edtext{eo rectae}{\lemma{eo}\Bfootnote{%
\textit{(1)}~ponder %
\textit{(2)}~lineae %
\textit{(3)}~rectae~\textit{L}}}
\textit{DF}, \textit{AB} licet inaequales sint aequiponderantes,\indexS{rectae aequiponderantes} 
quaeritur punctum \textit{G}. Ponatur inventum, et per ipsum ducatur recta \textit{PGQ} normalis 
\edtext{ipsi \textit{EGC}. Et quoniam}{\lemma{ipsi \textit{EGC}}\Bfootnote{%
\textit{(1)}~retentoque %
\textit{(2)}~. Et quoniam~\textit{L}}}
corpora \textit{DF} et \textit{AB} manent in aequilibrio\indexS{aequilibrium} 
quomodocunque circa centra gravitatis\indexS{centrum gravitatis} propria, \textit{E}, et \textit{C}, gyrentur, ideo ex infinitis 
\edtext{modis quibus invicem et ad libram\indexS{libra} inclinari possunt, eligatur}{\lemma{modis}\Bfootnote{%
\textit{(1)}~sola 
\textit{(2)}~eligat %
\textit{(3)}~quibus invicem et ad %
\textit{(a)}~l
\textit{(b)}~jugum %
\textit{(c)}~libram~\textit{L}}}
inclinatio hujusmodi, ut si in \textit{AB} sumatur \textit{BZ} 
\edtext{aequ. \textit{DF} et}{%
\lemma{aequ. \textit{DF}}%
\Bfootnote{%
\textit{(1)}~sit \textit{XY} %
\textit{(2)}~et~\textit{L}}}
hujus punctum medium \textit{X} ac residui \textit{AZ} medium \textit{R}. Sit \textit{XY} perpendicularis in \textit{PQ} (si opus productum) aequalis ipsi \textit{EG}, et punctum \textit{R} cadat in ipsam \textit{PQ} si opus productum. Hoc enim posito manifestum est omnia circa axem\indexS{axis} \textit{PQ} in aequilibrio\indexS{aequilibrium} fore, 
\edtext{nam \textit{BZ} et 
\lbrack\textit{DF}\rbrack\ 
aequales}{\lemma{nam \textit{BZ}}\Bfootnote{%
\textit{(1)}~ipsi \textit{DE} aequalis %
\textit{(2)}~et~\bars{\textit{DE} \textit{ändert Hrsg.}} aequales~\textit{L}}}
et ab axe (quoad medium punctum \textit{X} 
\edtext{vel \textit{E} ob}{\lemma{vel \textit{E}}\Bfootnote{%
\textit{(1)}~) aequidistantes %
\textit{(2)}~ob~\textit{L}}}
distantias \textit{EG} et \textit{XY} aequales) aequidistantes, utique in aequilibrio\indexS{aequilibrium} 
erunt et quae superest recta \textit{ARZ} quia in medio puncto \textit{R}
ab axe aequilibrii\indexS{axis aequilibrii} secatur, etiam in aequilibrio erit. 
\pend
\pstart
\href{\Scan{LH_37_03_082v}}%
{\lbrack82~v\textsuperscript{o}\rbrack}
Hunc ergo situm in quo id contingit investigare problema est pure geometricum adeoque in 
\edlabel{37_03_082-083_4a}%
\edtext{}{\xxref{37_03_082-083_4a}{37_03_082-083_4b}\lemma{potestate.}\Bfootnote{%
\textit{(1)}~Hoc ita sovemus.\lbrack\textit{!}\rbrack\ Quoniam \textit{XB} aequalis \textit{FE}, ergo dabitur et \textit{CX} %
\textit{(a)}~centro %
\textit{(b)}~quo radio ex \textit{C} describatur circulus. %
\textit{(aa)}~Et 
\textit{(aaa)}~pro
\textit{(bbb)}~producatur \textit{EG} in \textit{W}, ut fiat aequalis %
\textit{(bb)}~In \textit{EG} %
\textit{(cc)}~\bars{\textit{GC} \textit{streicht Hrsg.}} si opus producta sumatur \textit{GW} aequ. \textit{EG} \bars{unde \textit{streicht Hrsg.}} demissa \textit{WX} occurrat circulo centro \textit{C}. %
\textit{(aaa)}~de %
\textit{(bbb)}~radio \textit{CX} descripto in \textit{X}. %
\textit{(2)}~Hoc ita solvemus: %
\textit{(a)}~sit \textit{EG} aequ. \textit{YX} aequ. \textit{e}. et \textit{CE} aequ. \textit{c} erit \textit{CG} aequ. $c-e$. et \textit{CW} aequ. $GW-CG$ aequ. %
\textit{(aa)}~$e-$ %
\textit{(bb)}~$e-c+e$ aequ $2e-c$ %
\textit{(b)}~\textit{GC} est~\textit{L}}}%
potestate. 
\pend
\pstart
Hoc ita solvemus:
\textit{GC} est\edlabel{37_03_082-083_4b}
ad \textit{YX} (seu \textit{EG}) ut \textit{RC} ad \textit{RX}. Est autem \textit{RC} aequ. $AB-AR-BC$ 
\edtext{seu quia \textit{BC} aequ. \rule[-3mm]{0pt}{9mm}$\displaystyle\frac{AB}{2}$ fiet: \textit{RC} aequ. \rule[-3mm]{0pt}{9mm}$\displaystyle\frac{AB}{2}-AR$,}{%
\lemma{}\Bfootnote{seu quia \lbrack...\rbrack\ $-\,AR$, \textit{erg.~L}}}
et \textit{AR} 
\edtext{aequ. $\displaystyle\frac{AZ}{2}$ et \textit{AZ}}{\lemma{aequ. $\displaystyle\frac{AZ}{2}$}\Bfootnote{%
\textit{(1)}~et \textit{AZ} aequ. $AB\,-$
\textit{(a)}~\textit{DF}
\textit{(b)}~(\textit{ZB})\textit{DF} %
\textit{(2)}~et \textit{AZ}~\textit{L}}}
aequ. $AB-(ZB\ \text{vel})\; DF$ ergo \textit{RC} aequ. $\ovalbox{\displaystyle\frac{AB}{2},-\displaystyle\frac{AB}{2}}+\displaystyle\frac{DF}{2}$. At \textit{RX} aequ. $RC+CB-(XB\ \text{vel})\;\displaystyle\frac{DF}{2}$\rule[-3mm]{0pt}{9mm} seu \textit{RX} aequ. \rule[-3mm]{0pt}{9mm}$+\displaystyle\frac{DF}{2}+\displaystyle\frac{AB}{2}-\displaystyle\frac{DF}{2}$ 
\edtext{aequ. $\displaystyle\frac{AB}{2}$. Habemus}{\lemma{aequ. 
\rule[-1.8ex]{0pt}{5.4ex}%
$\displaystyle\frac{AB}{2}$,}\Bfootnote{%
\textit{(1)}~,~vel sic: \textit{CZ} aequ. 
\rule[-1.8ex]{0pt}{5.4ex}%
$\displaystyle\frac{AB}{2}-DF$ %
\textit{(2)}~.~Habemus~\textit{L}}}
ergo \textit{RC} aequ. \textit{DE} (vel \textit{XB}) et \textit{RX} aequ. \textit{CB}. Itaque \textit{CG} ad (\textit{XY} vel) \textit{GE} ut (\textit{DE} ad \textit{CB} vel ut) \textit{DF} ad \textit{AB}. Habemus ergo punctum \textit{G}. 
\pend
\pstart
Eo jam habito in recta \textit{GC} si opus producta sumatur \textit{GW} aequ. \textit{GE} et centro \textit{C} radio \textit{CX} describatur circulus inde ex \textit{W} demittatur parallela ipsi \textit{GQ} circulum secans in \textit{X}, rursusque 
\edtext{centro \textit{C} diametro}{\lemma{centro \textit{C}}\Bfootnote{%
\textit{(1)}~radio %
\textit{(2)}~diametro
\textit{(3)}~diametro~\textit{L}}}
\textit{CA} describatur circulus cui producta \textit{XC} occurrat in \textit{AB},
habemus \textit{AB} ita sitam, ut initio desiderabatur. Quoniam autem \textit{RX} aequ. \textit{AC}, hinc si sumatur \textit{CH} aequ. \textit{YX} vel \textit{GE} fiet triang. \textit{ACH} aequ. et simile Triangulo
\textit{RXY} et angulus \textit{AHC} et rectus et quia \textit{CH} aequ. \textit{EG} erit \textit{EH} aequ. \textit{CG}. 
\edtext{Porro $HC^2+AC^2$}{\lemma{}\Bfootnote{Porro \bars{aequ. \textit{streicht Hrsg.}} $HC^2$~\textit{L}}}
aequ. $\langle AH^2\rangle$ aequ. $-AC^2-EG^2$. 
\pend
\pstart
Est autem \textit{EG},\textit{DE} aequ. \textit{GC},\textit{AC} et $CG+GE$ aequ. \textit{c}.
Fiet $CG+\displaystyle\frac{GC,\!AC}{DE}$ aequ. \textit{c}. seu $\displaystyle\frac{c,\!DE}{AC+DE}$ aequ. \textit{CG} et $\displaystyle\frac{c,\!AC}{AC+DE}$ aequ. \textit{EG}.
Ergo $AH^2$ aequ. $AC^2-\displaystyle\frac{c^2,\!\overline{AC}^2}{AC^2+DE^2,\!2AC,\!DE}$.
Seu $AH^2$ aequ. $AC\!+\!EG,AC\!-\!EG$ id est $AC+\displaystyle\frac{c,\!AC}{AC+DE},\!+AC-\displaystyle\frac{c,\!AC}{AC+DE}$.
Unde non est necesse fieri angulum \textit{EAC} rectum seu $AH^2$ aequ. $\displaystyle\frac{c^2,\!DE^2}{AC+DE,\!\framebox{$2$}}$.%
Hoc tamen possumus efficere si lubet quia ipsius \textit{DF} quae sit inclinatio nil refert. Ergo possumus facere ipsi \textit{AB} perpendicularem vel etiam parallelam. \pend 
